\chapter*{\chapterstyle{Algèbre I {:} Corps des Complexes}}
On note \(\C\) l'extension du corps \(\R\) avec les deux lois usuelles. \< 

Est introduit le nombre imaginaire \(i\) dont le carré vaut \(-1\) et chaque nombre complexe se définit alors comme des sommes ou produits de réels et du nombre imaginaire.\+
On définit ainsi la \textbf{forme algébrique} d'un nombre complexe comme \(a + bi\) avec \(a, b \in \R\). \<

Soit \(z, z' \in \C\) et \(\theta \in \R/2\pi\). On notera \(\rho\) le module de \(z\) pour faciliter la lecture.

\subsection*{\subsecstyle{Module {:}}}
On appelle \textbf{module} de \(z\) le nombre réel positif \(|z|\) tel que {:}
\[
    \boxed{
        |z| = \sqrt{a^2 + b^2} = \sqrt{z\overline{z}} 
    }
\]
C'est un \textbf{prolongement} de la fonction valeur absolue à \(\C\).
On note que c'est une opération \textbf{multiplicative} mais \textbf{pas} additive, en effet on a l'inégalité triangulaire:
\[
   |z + z'| \leq |z| + |z'|   
\]
\subsection*{\subsecstyle{Forme trigonométrique {:}}}
On montre qu'il existe un angle \(\theta\) tel que:
\[
    \boxed{
       z = \rho(\cos\theta+i\sin\theta)   
    }
\]
Cette forme est appelée \textbf{forme trigonométrique} de \(z\) et elle est unique modulo \(2\pi\).\+
Elle se déduit du nombre \(\frac{z}{|z|} \in \U\) et de l'identité trigonométrique fondamentale.\<

On appelle \(\theta\) ainsi obtenu \textbf{l'argument} de \(z\).
\subsection*{\subsecstyle{Forme exponentielle {:}}}
De même pour \(\theta\) l'argument de \(z\), on a:
\[
    \boxed{
        z = \rho e^{i\theta} 
    }
\]
Cette forme est appelée \textbf{forme exponentielle} de \(z\) est également unique modulo \(2\pi\).\<

On montre alors les \textbf{formules d'Euler}:
\[
  \begin{aligned}
     \cos\theta := \frac{e^{i\theta}+e^{-i\theta}}{2}
  \end{aligned}
  \hspace*{40pt}
  \begin{aligned}
     \sin\theta := \frac{e^{i\theta}-e^{-i\theta}}{2i}
  \end{aligned}
\]
Ces formules se déduisent de (\(\ast\)). Elle sont très importantes car elle permettent de \textbf{linéariser} des expression trigonométriques.\<

Les propriétés de l'exponentielle nous permettent aussi d'obtenir:
\begin{align*}
   \arg(zz') &\eqmod{2\pi} \arg(z) + \arg(z') \\
   \arg(\frac{z}{z'}) &\eqmod{2\pi} \arg(z) - \arg(z')
\end{align*}

\subsection*{\subsecstyle{Conjugué {:}}}
On appelle \textbf{conjugaison} l'automorphisme involutif qui à \(z\) associe son \textbf{conjugué}, noté \(\overline{z}\) tel que:
\[
    \boxed{
        \overline{z} := a - bi = \rho\cos\theta - i\sin\theta = \rho e^{-i\theta}
    }
\]
C'est un morphisme \textbf{additif} et \textbf{multiplicatif}.

Graphiquement, le conjugué d'un nombre complexe est \textbf{le symétrique} de \(z\) par rapport à l'axe réel.\<

(\(\ast\)) On a alors la propriété suivante {:}
\[
  \begin{aligned}
     \Re(z) := \frac{z + \overline{z}}{2}
  \end{aligned}
  \hspace*{40pt}
  \begin{aligned}
     \Im(z) := \frac{z - \overline{z}}{2i}
  \end{aligned}
\]
\subsection*{\subsecstyle{Formule de Moivre {:}}}
Un propriété importante des formes trigonométriques et exponentielle est:
\[
   {(e^{i\theta})}^n = e^{n(i\theta)}
\]
De manière analogue, on a:
\[
   {(\cos\theta + i\sin\theta)}^n = \cos n\theta + i\sin n\theta
\]
Ici, on a choisi de considérer \(z \in \U\) mais ces propriétés sont vraies pour \textbf{tout nombre complexe}, il suffit alors de factoriser par le module pour se ramener au cas ci-dessus.
\subsection*{\subsecstyle{Racines n-ièmes {:}}}
Soit \(n \in \N\), une partie importante des problèmes impliquant des nombres complexes proviennent d'équations d'inconnue \(Z\) de la forme:
\[
   Z^n = z
\]
On peut montrer que l'ensemble des solutions de ce type de problème est:
\[
   \boxed{
      S := \Bigl\{ \sqrt[n]{\rho}e^{i\frac{\theta + 2k\pi}{n}}\; ; \; k \in \bigl\{0, 1, \ldots , n-1 \bigl\} \Bigl\}   
   }
\]
\underline{Cas particulier {:}}
Si on a une racine n-ième \(Z_0\) de \(Z\) et qu'on connaît les racines n-ièmes de l'unité, alors on peut obtenir toutes les racines n-ièmes de \(Z\) grâce à:

\[
   \Bigl\{ Z \in \C \; ; \; Z^n = z \Bigl\} = \Bigl\{ Z_0u \; ; \; u \in \U_n \Bigl\}   
\]
\subsection*{\subsecstyle{Le nombre complexe \(j\) {:}}}
On note \(j\) la première racine troisième de l'unité.
Le nombre \(j\) est singulier, car il vérifie:
\[
   j^2 = j^{-1} = \overline{j}
\]
Graphiquement, on peut observer que les images des nombres \(1\), \(j\) et \(\overline{j}\) forment un triangle équilatéral inscrit dans le cercle trigonométrique.

\chapter*{\chapterstyle{Algèbre I {:} Anneau des Polynômes}}
Soit K un corps, et \(n \in \N\), on appelle \textbf{polynômes} à coefficients dans K en l'indéterminée \(X\) les éléments de l'ensemble:
\[
   K[X] := \Biggl\{ \; \sum_{i=0}^n a_{i}X^{i} \; ; \; a_i \in K \;\Biggl\}
\]
On adjoint à cet ensemble deux lois + et \(\times\) qui \textbf{prolongent} les opérations correspondantes de K, pour obtenir \textbf{l'anneau des polynômes} \textbf{(K, +, \(\times\))}. Cet anneau est intègre car K est un corps, donc un anneau intègre.\<

Soit \(P, Q, R \in K[X]\).

\subsection*{\subsecstyle{Degré et Valuation {:}}}
On définit tout d'abord une propriété fondamentale appelée \textbf{degré} de \(P\) telle que deg\((P)\) est le plus grand coefficient non nul de P.
On a alors les propriétés du degré ci-dessous:
\begin{align*}
   \deg(P + Q) &\leq \max(\deg(P), \deg(Q))\\
   \deg(PQ) &= \deg(P) + \deg(Q)
\end{align*}
La valuation est définie de manière analogue comme le plus petit coefficient non nul de P.
\subsection*{\subsecstyle{Opérations {:}}}
On considère ci-dessous que deg\((P)\) = \(n\) et deg\((Q)\) = \(m\) et on note \(a_i, b_i\) les coefficients de \(P\) resp. \(Q\).\+
La lci d'addition est simplement effectuée terme à termes.\+
La lci de multiplication se définit explicitement comme suit:
\[
   PQ := \sum_{k=0}^{n+m}\sum_{i+j = k} a_{i}b_{i}X^{k}
\]

On ajoute aussi une opération appelée \textbf{dérivation formelle} d'un polynôme qui s'effectue comme la dérivation analytique usuelle.
\subsection*{\subsecstyle{Divisibilité {:}}}
On définit tout d'abord une \textbf{division euclidienne} de deux polynômes qui se comporte comme la division euclidienne usuelle, à la différence que la condition d'arrêt porte sur le \textbf{degré du reste} qui doit être inférieur à celui du diviseur.\<

Soit \(U, V \in R[X]^2\), on peut aussi définir une relation de \textbf{divisibilité} entre deux polynômes, cette relation est \textbf{transitive et réflexive} mais aussi \textbf{linéaire}, ie on a:
\[
    \boxed{
        D|A \land D|B \implies D| UA + VB   
    }
\]
\subsection*{\subsecstyle{Racines et Factorisation {:}}}
On dit que \(\alpha\) est une \textbf{racine} de \(P\) si \(P(\alpha) = 0\).\+
On montre alors le théorème fondamental ci-dessous:
\[
   \boxed{
      (X - \alpha) | P \Longleftrightarrow P(\alpha) = 0 
   }
\]
\pagebreak

On appelle \textbf{multiplicité d'une racine} \(\alpha\) l'entier \(m\) tel que:
\[
   \Bigr[ (X - \alpha)^m | P \Bigr] \; \land \; \Bigr[(X - \alpha)^{m+1} \nmid P\Bigr]
\]
Si on note \(P^{m}\) la dérivée n-ième de \(P\), on a aussi:
\[
   P^{m}(\alpha) = 0 \; \land \; P^{m + 1}(\alpha) \neq 0
\]
On en déduit que pour une racine \(\alpha\) de multiplicité \(m\), on peut \textbf{factoriser} \(P\) par \((X-\alpha)^m\).
Si on considère maintenant plusieurs racines \textbf{distinctes} \(a_0, a_1, \ldots, a_{n-1}, a_n\) de multiplicité respectivement \(m_0, m_1, \ldots, m_{n-1}, m_n\), on peut montrer la propriété suivante:
\begin{flalign*}
   \Bigr[ \prod_{i=0}^{n}(X-\alpha_i)^{m_i} \Bigr] \; \Biggr| \; P \shorteqnote{(On peut factoriser par le produit des \((X-\alpha_i)^{m_i}\))}
\end{flalign*}
\subsection*{\subsecstyle{Décomposition {:}}}
On appelle décomposition d'un polynôme \(P\) une factorisation \textbf{irréductible} de \(P\), cette décomposition dépend du corps considéré, en effet si on considère \(\C[X]\), on peut montrer le \textbf{théorème fondamental de l'Algèbre} ci-dessous:
\begin{center}
    \textbf{Tout polynôme de degré supérieur à 1 admet au moins une racine dans \(\C\)}.
\end{center}
Par suite, on peut montrer que tout les polynômes de \(\C\) sont \textbf{scindés}, ie que tout les polynômes irréductibles de \(\C[X]\) sont de degré 1.
Par contre dans \(\R[X]\), il existe des polynômes de degré 2 irréductibles (ceux dont le discriminant est négatif).\<

Soit \(z \in \C\), il existe une propriété très utile pour décomposer un polynôme réel qui est:
\[
   P(z) = 0 \implies P(\overline{z}) = 0
\]
\subsection*{\subsecstyle{Fonctions symétriques des racines {:}}}
On définit le \(k\)-ième \textbf{polynôme symétrique} à \(n\) indéterminées comme la somme des produits à \(k\) facteurs de ses indéterminées, et on le note \(\sigma_k\). \<

Par exemple pour un polynôme en quatre indéterminées \(a, b, c, d\), on a successivement:
\begin{flalign*}
    &\sigma_1 = a + b  + c + d \shorteqnote{(Somme des produits à 1 facteur)}\\
    &\sigma_2 = ab + ac + ad + bc + bc + cd \shorteqnote{(Somme des produits à 2 facteurs)}\\
    &\sigma_3 = abc + abd + acd + bcd \shorteqnote{(Somme des produits à 3 facteurs)}\\
    &\sigma_4 = abcd \shorteqnote{(Somme des produits à 4 facteurs)}
\end{flalign*}
De manière générale on a:
\[
    \sigma_k(x_1, \ldots, x_n) = \sum_{1 \leq i_1 < \ldots < i_k \leq n} x_{i_1}x_{i_2} \ldots x_{i_k}
\]
Soit \(P \in K[X]\) un polynôme de degré \(n\), on note \(x_1, x_2, \ldots, x_{n-1}, x_n\) ses \(n\) racines et \(a_1, a_2, \ldots, a_{n-1}, a_n\) ses \(n\) coefficients.\+
Alors, pour tout \(k \in \inticc{0}{n}\), on a le \textbf{théorème}:
\[
    \fbox{
        \sigma_k(x_1, x_2, \ldots, x_{n-1}, x_n) = (-1)^{k}\frac{a_{n-k}}{a_n}
    }
\]

Ce théorème permet d'obtenir une relation \textbf{coefficient-racines}.\+
Par exemple:
\begin{flalign*}
    &\bullet \text{Pour \(k = 1\), on trouve que la somme des racines de \(P\) vaut \(-\frac{a_n-1}{a_n}\)}\\
    &\bullet \text{Pour \(k = 2\), on trouve que la somme des doubles produits des racines de \(P\) vaut \(\frac{a_n-2}{a_n}\)}
\end{flalign*}



\chapter*{\chapterstyle{Algèbre I {:} Espaces Vectoriels}}
Soit \(E\) un ensemble \textbf{non-vide}, \(K\) un corps commutatif, et \(u \in E\), on dit alors que \((E, +, \cdot)\) est un \textbf{espace vectoriel sur K} si les conditions ci-dessous sont réunies:
\begin{align*}
   &\bullet \;\;\; (E, +) \text{ est un \textbf{groupe abélien}} \\
   &\bullet \;\; \text{ La loi \(\cdot\) est une application de \(K \times E \longrightarrow E\)} \\
   &\bullet \;\; \text{ La loi \(\cdot\) vérifie \textbf{la distributivité mixte et l'associativité mixte}} \\
   &\bullet \;\;\; 1 \cdot u = u
\end{align*}

On appelle alors \textbf{vecteurs} les éléments de \(E\) et \textbf{scalaires} les éléments de \(K\).

\subsection*{\subsecstyle{Sous-espaces vectoriels {:}}}
Soit \(F \subseteq E\), on dit que \(F\) est un \textbf{sous-espace vectoriel} ssi:
\begin{align*}
   &\bullet \;\; F \text{ est non vide.} \\
   &\bullet \;\; F \text{ est stable par somme.} \\
   &\bullet \;\; F \text{ est stable par multiplication externe.}
\end{align*}

Il est utile de noter que \textbf{l'intersection} de deux sous-espaces vectoriels est aussi un sous-espace vectoriel et que par contre, l'union de deux sous-espaces vectoriels n'est \textbf{en général} pas un sous-espace vectoriel.\<

Un cas particulier est celui du \textbf{sous-espace vectoriel engendré} par \(F\) qu'on note:
\[
   \text{Vect(\(F\))} := \Biggl\{ \sum_{i=0}^{n} \lambda_i u_i\; ; \; \lambda_{i} \in K \; , \; u_{i} \in E \Biggl\}
\]
Plus concrètement, Vect(\(F\)) est \textbf{l'ensemble des combinaisons linéaires} des vecteurs de \(F\).\<

Une propriété fondamentale est que Vect(\(F\)) est un \textbf{opérateur de clôture} par combinaison linéaire, ie c'est une application \textbf{idempotente, croissante et extensive}.


\subsection*{\subsecstyle{Familles libre et génératrices {:}}}
On appelle \textbf{famille} de vecteurs un n-uplet de la forme \((u_1, u_2, \ldots, u_{n-1}, u_n)\).\+
Soit \(\Fam := (u_1, u_2, \ldots, u_{n-1}, u_n)\) une famille de vecteurs de \(E\).\<

On dit que \(\Fam\) est \textbf{génératrice} de E si on a \(E = \text{Vect(\(\Fam\))}\).\+
On dit que \(\Fam\) est \textbf{libre} si tout combinaison linéaire nulle de vecteurs de \(\Fam\) est nécessairement à coefficient tous nuls.\+
Formellement, \(\Fam\) est libre si et seulement si on a:
\[
    \forall (\lambda_1, \lambda_2, \ldots, \lambda_{n-1}, \lambda_n) \in K^n \; ; \; \Biggr[ \sum_{i=0}^n \lambda_i u_i = 0_E \implies (\lambda_1, \lambda_2, \ldots, \lambda_{n-1}, \lambda_n) = (0, 0, \ldots, 0, 0) \Biggr]  
\]

Soit \(\Fam := (u_1, u_2, \ldots, u_k, \ldots, u_{n-1}, u_n)\) une famille libre et génératrice.\<

Si un vecteur \(v\) n'appartient pas à \(\text{Vect(\(\Fam\))}\), alors la famille étendue \((u_1, u_2, \ldots, u_{n-1}, u_n, v)\) est toujours libre.\+
Si un vecteur \(u_k\) appartient à \(\text{Vect(\(\Fam\))}\), alors la famille réduite \((u_1, \ldots, u_{k-1}, u_{k+1}, \ldots, u_n)\) est toujours génératrice.
\pagebreak

\subsection*{\subsecstyle{Bases {:}}}
On appelle \textbf{base} de \(E\) une famille \textbf{libre et génératrice}, et on a le théorème suivant:
\begin{center}
    \fbox{
        \textbf{Tout espace vectoriel possède une base.}
    }
\end{center}

Soit \(\mathscr{B} = (e_1, e_2, \ldots, e_{n-1}, e_{n})\) une base de \(E\).\+
On appelle \textbf{coordonnées} de \(u\) dans la base \(\mathscr{B}\) les coefficients de la décomposition de \(u\) dans la base \(\mathscr{B}\) et on note:
\[
    [u]^{\mathscr{B}} = [\lambda_1 e_1 + \lambda_2 e_2 + \ldots + \lambda_{n-1} e_{n-1} + \lambda_{n-1} e_{n-1} ]^{\mathscr{B}} = 
    \begin{pmatrix}
        \lambda_1\\
        \vdots\\
        \lambda_{n}
    \end{pmatrix}
\]

\subsection*{\subsecstyle{Sommes et espaces supplémentaires{:}}}
Soit \(F\), \(G\) et \(H\) trois sous-espaces vectoriels de \(E\).\+
On appelle \textbf{somme} de deux sous-espaces vectoriels l'ensemble:
\[
    F + G := \Bigl\{ u_1 + u_2 \; ; \; u_1 \in E \; , \; u_2 \in F\Bigl\}
\]

Si il y a \textbf{unicité} de la décomposition, on dit alors que \(F\) et \(G\) sont en \textbf{somme directe} et on note:
\[
    F \oplus G := \Bigl\{ u_1 + u_2 \; ; \; u_1 \in E \; , \; u_2 \in F \; , \; F \cap G = \{0_E\}\Bigl\}
\]

Si \(H\) se décompose de manière unique en somme de vecteurs de \(F\) et de \(G\), alors on dit que \(F\) et \(G\) sont \textbf{supplémentaires} dans \(H\) et on note:
\[
    \boxed{
        H = F \oplus G
    }
\]

Si on connaît une base \(\mathscr{B}_1 = (e_1, e_2, \ldots, e_{n-1}, e_{n})\) de F et une base \(\mathscr{B}_2 = (f_1, f_2, \ldots, f_{n-1}, f_{n})\) de G, et si on pose \(\Fam := (e_1, \ldots, e_{n}, f_1, \ldots, f_{n})\) on a aussi le théorème suivant:
\[
    H = F \oplus G \Longleftrightarrow \Fam \text{ est une base de H}
\]

\subsection*{\subsecstyle{Hyperplans {:}}}
Soit \(F\) un sous-espace vectoriel de \(E\).\+
\(F\) est un \textbf{hyperplan} de E si et seulement si il vérifie l'une des conditions équivalentes suivantes:
\begin{align*}
    &\bullet \text{ Il existe \textbf{une droite vectorielle} \(G\) de \(E\) telle que \(F \oplus G = E\).} \\
    &\bullet \text{ \(H\) est défini par une \textbf{équation linéaire homogène} non triviale.}
\end{align*}
Informellement, les hyperplans sont les \textbf{sous-espaces propres maximaux} de \(E\). \<

On remarque alors que résoudre un système d'équations homogène revient à chercher \textbf{l'intersection} des hyperplans qui sont décrits par ces équations.\<

En d'autres termes, l'ensemble des solutions d'un système d'équations homogène est nécessairement lui aussi \textbf{un sous-espace vectoriel}, en tant qu'intersection de sous-espaces vectoriels.


